\documentclass[11pt, letterpaper]{article}
\usepackage[margin=1in]{geometry}
\usepackage{titlesec}
\usepackage{enumitem}
\usepackage{hyperref}
\usepackage{titling}
\usepackage{cite}

% Formatting section titles for a compact 1-page look
\titleformat{\section}{\large\bfseries}{\thesection}{1em}{}
\titlespacing*{\section}{0pt}{1.2ex plus 1ex minus .2ex}{0.5ex plus .2ex}
\setlength{\droptitle}{-8em} 
\title{\textbf{Project-X Proposal: Forcing Visual Grounding}\\
\large Causal-Aware Few-Shot Learning in Vision-Language Models via Counterfactual Retrieval and Attention Modulation}
\author{Simone Cama \\ 2044010}
% \date{May 10, 2026}

% remove spacing around date:
\predate{}
\postdate{}
\date{} 

\begin{document}

\maketitle
\vspace{-1cm} % Reduce space below title to ensure it fits on one page

\section*{Abstract}
Vision-Language Models (VLMs) have demonstrated remarkable capabilities in In-Context Learning (ICL), allowing them to adapt to new tasks using a few interleaved image-text demonstrations. However, recent literature exposes critical flaws in this paradigm: VLMs frequently suffer from modality imbalance (ignoring visual inputs in favor of textual bias) and rely on spurious correlations gathered from passive, similarity-based retrieval. This project proposes a pipeline designed to force causal visual grounding, introducing a counterfactual retrieval algorithm coupled with a custom cross-modal attention-forcing penalty, compelling the model to utilize visual evidence.

\section{Baseline \& Identified Limitations}
The chosen baseline is \textbf{OpenFlamingo} \cite{awadalla2023openflamingo}, an open-source replication of the foundational Flamingo architecture \cite{flamingo}. Based on a rigorous literature review, three primary methodological limitations have been identified:
\begin{itemize}[noitemsep, topsep=2pt]
    \item \textbf{Modality Imbalance (Visual Ignorance):} Research indicates that ICL in VLMs is predominantly driven by textual demonstrations, with visual information having a negligible effect due to suboptimal cross-attention mechanisms \cite{chen2024understanding}.
    \item \textbf{Spurious Correlations in Retrieval:} Standard Retrieval-based In-Context Example Selection (RICES) selects demonstrations based on passive similarity. This retrieves visually correlated but non-causal examples, preventing the model from learning underlying causal relationships \cite{xiong2026retrieving}.
    \item \textbf{The Many-Shot Paradox:} Attempting to resolve poor reasoning by scaling up the number of in-context examples actually degrades functional correctness, with optimal performance constrained to a narrow few-shot window \cite{oskooei2025many}.
\end{itemize}

\section{Proposed Ad-Hoc Methodological Novelty}
To overcome these limitations I propose a dual methodological novelty integrating both algorithmic and architectural optimizations:
\begin{enumerate}[noitemsep, topsep=2pt]
    \item \textbf{Algorithmic Novelty (Causal-Mixed Retrieval):} I will introduce a dynamic demonstration selection algorithm merging counterfactual Composed Image Retrieval (CIR) \cite{xiong2026retrieving} with mixed-modality semantic constraints \cite{chen2024understanding}. To avoid the Many-Shot Paradox \cite{oskooei2025many}, the retrieval context window will be strictly bounded dynamically (5--25 shots). 
    \item \textbf{Architectural Novelty (Cross-Modal Attention-Forcing Penalty):} To ensure the VLM utilizes the retrieved counterfactual images, I will implement a custom auxiliary loss function during a brief parameter-efficient fine-tuning phase. This custom loss will penalize the model mathematically if the attention weights assigned to \textit{counterfactual visual tokens} fall below a defined threshold relative to textual tokens, forcing the model to ground its reasoning in visual causal differences rather than textual hallucination.
\end{enumerate}

\section{Dataset Selection \& Implementation Pipeline}
Because this project evaluates \textit{In-Context Learning} (ICL), the model cannot be evaluated on data it has been fine-tuned on. To ensure rigorous validation and an apples-to-apples comparison with the foundational literature, the pipeline is divided into a general training phase and a strict zero/few-shot evaluation phase.

\textbf{1. Parameter-Efficient Fine-Tuning (PEFT) Phase:}\\
The \textbf{Cross-Modal Attention-Forcing Penalty} will be learned during a brief PEFT phase utilizing a held-out subset of the standard \textbf{VQAv2} dataset \cite{goyal2017making}. This ensures the model learns the mechanism for visual grounding on general data without contaminating the specialized downstream evaluation benchmarks.

\textbf{2. Zero-Shot \& Few-Shot Evaluation Phase:}\\
To rigorously test the novelties, the model will be evaluated on the following benchmarks explicitly utilized in the baseline literature:
\begin{itemize}[noitemsep, topsep=2pt]
    \item \textbf{Animals with Attributes 2 (AwA2):}
    \begin{itemize}[noitemsep]
        \item \textit{Description:} A zero-shot learning dataset containing 50 common animal classes annotated with 85 continuous visual attributes (e.g., ``has stripes'', ``long neck'', ``furry'').
        \item \textit{Purpose:} AwA2 validates the \textbf{Causal-Mixed Retrieval} algorithm by querying the model on specific counterfactual visual, mathematically proving whether the model successfully grounded its reasoning in the retrieved causal examples.
        \item \textit{Source/Link:} Hugging Face (\href{https://huggingface.co/datasets/quanda-bench-test/awa2}{\texttt{quanda-bench-test/awa2}}) \cite{xian2018zero}.
    \end{itemize}

    \item \textbf{GQA (Visual Reasoning in the Real World):}
    \begin{itemize}[noitemsep]
        \item \textit{Description:} A compositional Visual Question Answering (VQA) dataset explicitly engineered to minimize language biases and test real-world spatial relationships.
        \item \textit{Purpose:} GQA serves as the optimal benchmark to validate the \textbf{Cross-Modal Attention-Forcing Penalty} \cite{chen2024understanding}. Because the questions are heavily compositional (e.g., ``Is the man to the right or left of the dark cup?''), the VLM cannot rely on spurious correlations. If the proposed architectural penalty successfully forces visual grounding, accuracy on GQA's spatial subsets will directly improve.
        \item \textit{Source/Link:} Hugging Face (\href{https://huggingface.co/datasets/Mineru/GQA}{\texttt{Mineru/GQA}}) \cite{hudson2019gqa}.
    \end{itemize}
\end{itemize}


\section{Evaluation \& Ablation Study}
To isolate the impact of the ad-hoc novelties, a strict ablation study will be conducted:
\begin{itemize}[noitemsep, topsep=2pt]
    \item \textbf{Baseline:} OpenFlamingo + Standard RICES retrieval.
    \item \textbf{Ablation 1:} OpenFlamingo + Causal-Mixed Retrieval only (isolating the impact of the algorithm).
    \item \textbf{Ablation 2:} OpenFlamingo + Causal-Mixed Retrieval + Attention-Forcing Penalty.
\end{itemize}




\vspace{0.5cm}
% Bibliography
\bibliographystyle{plain}
\bibliography{references} % Points to references.bib

\end{document}